% META: {"id":"1NSI-TC-EX-022-CDP","chapitre":"1NSI-TYPES-CONSTRUITS","type_objet":"coup_de_pouce","exercice_ref":"1NSI-TC-EX-022","status":"approved","origin":"needs_review","status_history":["needs_review","audited","accepted","approved"],"release_acceptance":"RELEASE_OWNER_BATCH_ACCEPTANCE","acceptance_closure_digest":"sha256:9b3ccf9a81c5520fbb7e03b7b2d3e7bf2057a02834b6d908bf3be5f49f3b553f","acceptance_receipt_digest":"sha256:4fad86f0282e0fa4e0419cca1ad6379ae7af5a280c04a4a90880f90a1c044242"}
\coupDePouce{1}{On te demande de tracer l'exécution d'un programme qui utilise deux tableaux parallèles (élèves et notes) et affiche ceux qui ont la moyenne.}
\coupDePouce{2}{L'indice \lstinline{i} sert à accéder à la même position dans les deux tableaux : \lstinline{eleves[i]} et \lstinline{notes[i]} vont ensemble. C'est pour cela qu'on utilise \lstinline{range(len(eleves))} et non \lstinline{for e in eleves}.}
